% ============================================================
%  SmartGrid AI Audit Framework — DJES Journal Paper Body
%  Body only: \section{Introduction} through \section{Conclusion}
%  No preamble. No \begin{document}.
% ============================================================

% ------------------------------------------------------------------
\section{Introduction}
\label{sec:intro}
% ------------------------------------------------------------------

Modern electric power networks have changed significantly over the past two decades. What were once
passive, one-way distribution systems have evolved into two-way cyber-physical networks that carry
both power and information simultaneously. Sensors, intelligent electronic devices, phasor
measurement units (PMUs), and advanced metering infrastructure now produce continuous telemetry
streams from thousands of assets. This operational intelligence lets grid operators detect faults
faster and balance loads more efficiently, but it also expands the attack surface that adversaries
can exploit.

Three classes of cyber attack are particularly damaging in smart-grid settings. False data injection
(FDI) attacks corrupt sensor measurements at the telemetry layer. An adversary with knowledge of
the state-estimation equations can craft measurement vectors that pass conventional bad-data
detection tests while silently biasing the grid model used for control decisions. A 15\% positive
bias on generator voltage readings, sustained over several hours, can cause the energy management
system to underestimate generation capacity and trigger unnecessary load shedding. Denial-of-service
(DoS) attacks flood or block communication channels. When PMU telemetry is blacked out for even a
30-second window, state estimation degrades and protection relays lose the situational awareness
needed to distinguish a fault from a switching event. Man-in-the-middle (MITM) attacks intercept
and modify messages in transit. A field-device packet whose integrity field drops by 35\% while its
payload jumps outside the physical rate-of-change limit is a strong indicator that a relay or
gateway has been compromised.

Beyond detection, there is a resource-allocation problem that receives less attention in the
literature. A utility operating 100 or more intelligent assets cannot schedule full audits on every
device at every time step; the operational cost would be prohibitive. Yet a static round-robin
schedule gives attackers a predictable window between audits. A security-aware scheduler must
allocate audit effort dynamically, concentrating resources on high-risk assets while maintaining
baseline coverage everywhere.

Prior work has addressed pieces of this problem but not the whole. Statistical deviation methods
such as CUSUM and EWMA are computationally light but cannot distinguish physical noise from
deliberate injection. Machine learning methods including Isolation Forest and One-Class SVM improve
discrimination but have high false-positive rates when applied to heterogeneous agent populations
without per-type profiles. Deep learning approaches based on LSTM networks capture temporal
dependencies but, when tuned for low false-positive rates, produce recall values below 25\%. None
of these prior systems couples detection to an adaptive audit scheduler, and none has been
integrated with a live supervisory control and data acquisition (SCADA) platform. Explainability ---
the ability to tell an operator which specific features drove a detection decision --- has received
almost no attention in the grid-security domain despite its importance for regulatory compliance and
operator trust.

This paper presents the SmartGrid AI Audit Framework, an end-to-end prototype that addresses all of
the above gaps. The system was evaluated on a population of 100 agents spanning four types
(generators, substations, PMUs, and circuit breakers) over a 24-hour simulation cycle comprising
28,800 time steps with three injected attack scenarios. The main contributions are:

\begin{enumerate}
    \item A three-layer anomaly detection pipeline that combines Mahalanobis deviation scoring, a
    dual-branch LSTM network, and attack-type-specific sub-detectors (CUSUM for FDI, a 2-of-3
    network-rule detector for DoS, and an integrity-plus-temporal-jump detector for MITM), unified
    by an OR-with-precedence combiner that suppresses false-positive inflation.

    \item A hybrid audit scheduler that pairs tabular Q-learning (state space: 12 states, action
    space: 3 escalation levels) with a gradient-descent frequency refinement step, producing
    continuous per-agent audit frequencies that respect a total budget constraint of 150 units per
    cycle.

    \item A feature-level explainability module that ranks the top-5 contributing features for every
    detection decision, refreshed at each detection cycle, providing operators with actionable
    context rather than binary alerts.

    \item A live integration layer connecting the detection and scheduling pipeline to Rapid SCADA v6
    via a PowerShell bridge, with a Next.js dashboard that presents all detection, scheduling, and
    audit data across eight distinct operational views.
\end{enumerate}

The remainder of this paper is organized as follows. Section~\ref{sec:lit} surveys related work.
Section~\ref{sec:methods} describes the threat model, system architecture, detection layers,
scheduler, explainability module, dashboard, experimental setup, and implementation stack.
Section~\ref{sec:results} reports results across nine sub-studies. Section~\ref{sec:conclusion}
states the conclusions, limitations, and future directions.

% ------------------------------------------------------------------
\section{Literature Review}
\label{sec:lit}
% ------------------------------------------------------------------

\subsection{Statistical and Rule-Based Detection}

The earliest automated anomaly detectors for industrial control systems applied classical
statistical process control methods. The CUSUM algorithm, introduced by Page~\cite{page1954} for
manufacturing quality control, detects a persistent shift in a monitored variable by accumulating
the signed deviation from a target mean. When the cumulative sum exceeds a threshold $h$, an alarm
is raised. For smart-grid telemetry, CUSUM is well suited to FDI because injection typically
produces a non-zero-mean bias in physical residuals; random sensor noise is zero-mean and the
statistic stays near zero in the absence of attack.

Chi-squared and generalized likelihood ratio tests have been applied to the state-estimation
residuals of power systems to detect bad data~\cite{liu2011false}. These methods work well when
the measurement model is well specified, but they require knowledge of the network admittance matrix
and are vulnerable to the stealthy FDI construction that Liu et al.\ originally identified: an
adversary who knows the topology can craft a vector $\mathbf{a}$ such that the injected measurements
$\mathbf{z}' = \mathbf{z} + \mathbf{a}$ produce the same residual as the unattacked case.

Exponentially weighted moving average (EWMA) charts provide a middle ground between sensitivity and
noise rejection by down-weighting old observations. Zhang et al.~\cite{zhang2019} applied EWMA
to voltage and frequency telemetry from substation automation systems and reported FPR below 2\%
under normal operating conditions. However, EWMA treats each sensor independently and does not
model cross-feature correlations; a coordinated attack that keeps individual readings within bounds
while shifting the correlation structure will evade it.

\subsection{Classical Machine Learning Methods}

Support vector machines (SVMs) were among the first machine learning models applied to grid
intrusion detection. The one-class variant trains a boundary around normal operating data in a
kernel-projected feature space and flags observations that fall outside the boundary. While
one-class SVM achieves high recall on isolated anomalies, its decision boundary is sensitive to the
choice of kernel and the fraction of training data treated as support vectors. On heterogeneous
populations where generators, substations, and PMUs have different statistical profiles, a single
shared model tends to produce high false-positive rates. The five-method comparison in the present
study (Section~\ref{sec:results}) finds one-class SVM accuracy of 46.35\% and FPR of 59.99\% on
the 100-agent dataset, confirming these known weaknesses.

Isolation Forest~\cite{liu2008isolation} builds an ensemble of random partition trees and flags
observations that are isolated by fewer splits than average. It is parameter-light and handles
high-dimensional data without explicit distance metrics. Several studies have applied Isolation
Forest to SCADA logs~\cite{beaver2013evaluation} with recall around 54\% and FPR near 28\%, figures
consistent with the 53.92\% recall and 27.95\% FPR observed in the present study. The core
limitation is that Isolation Forest has no mechanism to use domain knowledge about what FDI, DoS,
and MITM look like at the signal level; all anomaly evidence is treated as equivalent.

Random forests and gradient-boosted trees have been applied to grid anomaly detection with good
accuracy when labeled training data are available. Ashrafuzzaman et al.~\cite{ashrafuzzaman2020}
report random-forest accuracy above 95\% on a labeled dataset of SCADA intrusions, but note that
the class imbalance between normal and attack samples requires careful balancing and that the method
does not generalize well to attack variants not present in the training set.

\subsection{Deep Learning for Grid Anomaly Detection}

Recurrent neural networks, and LSTM networks in particular, have become the most widely studied
deep learning architecture for sequential SCADA anomaly detection. Kravchik and Shabtai~\cite{kravchik2018}
trained a single-input LSTM on the SWaT water treatment dataset and
reported F1 scores above 0.97. However, that dataset has a single process type; smart-grid agents
present multiple concurrent time series with different physical units and different noise
characteristics. Merging all features into a single input branch conflates physical and cyber
telemetry whose statistical behaviors differ substantially.

Several subsequent papers proposed multi-stream architectures that process physical and cyber
features separately before merging at a higher layer. CNN-LSTM hybrids add a convolutional front-end
to extract local temporal patterns before the recurrent layer processes the sequence. Khoei
et al.~\cite{khoei2021} applied a CNN-LSTM to a simulated smart-grid dataset and found that the
convolutional stage improved detection of short-duration spike attacks but offered no advantage for
sustained low-amplitude attacks that span many time steps. The present work uses a dual-branch LSTM
(one branch per feature domain) rather than a CNN-LSTM, based on the observation that spatial
convolution over temporal windows requires more training data than is available from a 24-hour
simulation.

A persistent challenge in deep learning for anomaly detection is the tradeoff between precision and
recall. Models tuned for low false-positive rate tend to produce recall below 25\% on stealthy
attacks. The multi-layer architecture in this paper addresses this by introducing a temporal
accumulator (Layer B) and attack-type sub-detectors (Layer C) that complement the primary LSTM
threshold path.

\subsection{Reinforcement Learning for Audit Scheduling}

The use of reinforcement learning (RL) for cyber-security resource allocation has grown since
Malialis and Kudenko~\cite{malialis2014} demonstrated that a multi-agent Q-learning system could
learn adaptive honeypot placement. In the smart-grid context, the audit scheduling problem is a
constrained resource allocation task: given a budget of $B$ audit units and $N$ agents with
time-varying risk scores, decide which agents to audit and with what frequency.

Priyadarsini et al.~\cite{priyadarsini2025} proposed a Q-learning scheduler for the grid-audit
problem with a discrete action space (increase, decrease, hold) and a reward function that penalizes
missed attacks. Their results show that dynamic scheduling reduces attack residence time compared to
fixed-interval scheduling. However, the Q-learning component alone produces discrete frequency steps
that do not optimize the continuous cost function. The hybrid approach in the present work adds a
gradient-descent refinement step that adjusts each agent's audit frequency along the gradient of a
cost function $C_i(f_i) = C_a f_i + C_f (R_i / f_i)$, converging to the analytically optimal
frequency $f_i^* = \sqrt{C_f R_i / C_a}$.

Deep Q-networks and policy-gradient methods have been applied to network intrusion response, but
the tabular Q-learning approach used here is sufficient for the 12-state, 3-action space and is
fully interpretable: the Q-table can be inspected directly to understand what the system has learned
about each risk category.

\subsection{Explainability in Security Systems}

The demand for explainable AI (XAI) in security-critical applications has grown as regulators in
the energy sector require documented justification for automated decisions. NERC CIP standards and
IEC 62443 both require audit trails and human-readable rationale for access and control decisions.

SHAP (SHapley Additive exPlanations) and LIME (Local Interpretable Model-agnostic Explanations)
have been applied to intrusion detection classifiers to produce per-feature importance
scores~\cite{arrieta2020explainable}. These methods are model-agnostic and produce reliable
attributions for tree-based classifiers, but they require additional computational passes and are
not always well suited to real-time streaming contexts. The present work uses a direct contribution
decomposition: for each flagged agent, the squared normalized deviation of each feature is computed
and expressed as a percentage of the total anomaly score. This is computationally inexpensive,
directly tied to the detection formula, and produces rankings that operators in grid-security
contexts find intuitive.

\subsection{Research Gaps Summary}

Table~\ref{tab:litgap} summarizes eight representative related works against the capabilities
addressed in this paper. Four gaps stand out consistently: (1) absence of live SCADA integration,
(2) no dynamic audit scheduling, (3) no explainability, and (4) evaluation on a single attack type
or on homogeneous agent populations. The SmartGrid AI Audit Framework addresses all four.

\begin{table*}[htbp]
\centering
\caption{Summary of representative related works and research gaps.}
\label{tab:litgap}
\small
\begin{tabular}{@{}llllllll@{}}
\hline
\textbf{Reference} & \textbf{Approach} & \textbf{Attack Types} & \textbf{Live Deploy} &
\textbf{Scheduling} & \textbf{Explain.} & \textbf{Multi-Agent} & \textbf{Key Gap} \\
\hline
Liu et al.\ (2011)~\cite{liu2011false} & Chi-squared, GLR & FDI only & No & No & No & No & Single attack type, no scheduling \\
Zhang et al.\ (2019)~\cite{zhang2019} & EWMA rule-based & DoS/FDI & No & No & No & No & No temporal modeling, no ML \\
Kravchik \& Shabtai (2018)~\cite{kravchik2018} & LSTM & General & No & No & No & No & No domain-specific attack layers \\
Khoei et al.\ (2021)~\cite{khoei2021} & CNN-LSTM & FDI, DoS & No & No & No & No & No scheduling or explainability \\
Ashrafuzzaman et al.\ (2020)~\cite{ashrafuzzaman2020} & Random Forest & Intrusions & No & No & Partial & No & Requires labeled data, no scheduler \\
Malialis \& Kudenko (2014)~\cite{malialis2014} & Multi-agent RL & DDoS & No & Yes (resource) & No & Yes & No power-grid specifics \\
Beaver et al.\ (2013)~\cite{beaver2013evaluation} & Isolation Forest & SCADA logs & Partial & No & No & No & High FPR, no attack type labeling \\
Priyadarsini et al.\ (2025)~\cite{priyadarsini2025} & Deviation + Q-learning & General & No & Yes & No & Yes & Single-layer, no SCADA, no XAI \\
\hline
\textbf{This work} & Multi-layer + RL + XAI & FDI, DoS, MITM & \textbf{Yes} & \textbf{Yes} & \textbf{Yes} & \textbf{Yes} & --- \\
\hline
\end{tabular}
\end{table*}

% ------------------------------------------------------------------
\section{Materials and Methods}
\label{sec:methods}
% ------------------------------------------------------------------

\subsection{Threat Model}
\label{sec:threatmodel}

The system is modeled as a population of $N = 100$ autonomous grid agents. Each agent represents a
physical asset that produces telemetry at each discrete time step. The adversary is assumed to have
partial insider knowledge: they know the agent identifiers and approximate operating baselines but
do not know the current detection thresholds or the Q-table state. This corresponds to a realistic
scenario where an attacker has conducted reconnaissance (e.g., by passively monitoring unencrypted
Modbus or DNP3 traffic) but has not compromised the security back-end.

Three attack vectors are modeled:

\textbf{False Data Injection (FDI).} The adversary injects a positive voltage bias of $+15\%$ on
the physical readings of generator agents (GEN-01 through GEN-20) at four-hour intervals. The
injection is designed to remain below the single-step detection threshold of any individual sensor
while accumulating a statistically significant bias over a multi-step window. In the 24-hour
simulation, FDI events occur at hours 0, 4, 8, 12, 16, and 20, for a total of 24 injected steps
per generator per event.

\textbf{Denial of Service (DoS).} Five PMU agents are selected at random from the pool
PMU-51 through PMU-75 at each injection event. Their cyber telemetry (latency, packet loss,
communication frequency) is set to a blackout profile for a window of 30 consecutive time steps.
The blackout window simulates a flooding attack that saturates the communication channel and
prevents any measurement from reaching the backend.

\textbf{Man-in-the-Middle (MITM).} Substation agents (SUB-21 through SUB-50) are subject to
integrity degradation and temporal jumps. The injected values have an integrity score reduced by at
least 35\% from the baseline and a physical-feature delta that exceeds the expected rate-of-change
limit. MITM events are injected at irregular intervals to simulate an attacker who avoids periodic
patterns that might be detected by frequency analysis.

The agent vulnerability map is shown in Table~\ref{tab:agenttypes}. Generators are the primary
target for FDI because their voltage and power readings feed directly into economic dispatch
calculations. PMUs are the primary DoS target because loss of PMU data has an outsized effect on
state-estimation accuracy. Substations are the MITM target because their gateway function makes
them the natural interception point.

\subsection{System Architecture}
\label{sec:arch}

The framework operates across two loosely coupled workspaces that share a common backend API.

The \textit{Simulation Workspace} generates synthetic telemetry internally, injects attacks
according to the schedule defined in Section~\ref{sec:threatmodel}, and streams the resulting
feature vectors to the detection pipeline at each time step. This workspace is used for the
controlled experiments in Section~\ref{sec:results}.

The \textit{Live SCADA Workspace} connects to a Rapid SCADA v6 installation running 100 calculated
channels. A PowerShell bridge polls the Rapid SCADA Web API at five-second intervals, packages the
channel readings into JSON batches, and forwards them to the FastAPI backend via the
\texttt{/v1/scada/ingest/tags/batch} endpoint. The backend normalizes the incoming tags against
per-agent profiles and feeds the resulting feature vectors into the same detection pipeline used in
the simulation workspace.

The overall data flow has five tiers:

\begin{enumerate}
    \item \textbf{Data tier.} Rapid SCADA (live) or the simulation engine generates raw telemetry.
    \item \textbf{Ingestion tier.} The PowerShell bridge or the simulation loop delivers batches to
    the FastAPI backend.
    \item \textbf{Detection tier.} The multi-layer pipeline scores each agent and produces a binary
    flag, an attack-type label, a confidence score, and a feature attribution vector.
    \item \textbf{Scheduling tier.} The hybrid Q-learning and gradient-descent scheduler reads the
    current risk vector and outputs per-agent audit frequencies within the budget constraint.
    \item \textbf{Presentation tier.} The Next.js dashboard retrieves data via REST polling and
    WebSocket push and renders it across eight views.
\end{enumerate}

\subsection{Agent Types and Feature Specification}
\label{sec:agents}

Table~\ref{tab:agenttypes} lists the four agent types, their counts, channel assignments, key
features, and primary threat vectors. Each agent contributes three calculated channels to Rapid
SCADA, giving 300 total channels for the 100-agent population. Features are divided into physical
measurements and cyber-layer metrics, enabling the dual-branch LSTM to process each domain with
architecturally appropriate parameters.

\begin{table*}[htbp]
\centering
\caption{Agent type specification: counts, channel ranges, features, and primary threats.}
\label{tab:agenttypes}
\small
\begin{tabular}{@{}llllll@{}}
\hline
\textbf{Type} & \textbf{IDs} & \textbf{Count} & \textbf{Key Physical Features} &
\textbf{Key Cyber Features} & \textbf{Primary Threat} \\
\hline
Generator (GEN) & GEN-01 -- GEN-20 & 20 & Voltage (230 V), Current (15 A), Power (3 kW) &
Latency (3 ms), Packet loss (0.001), Integrity (1.0) & FDI (voltage bias) \\
Substation (SUB) & SUB-21 -- SUB-50 & 30 & Voltage (230 V), Power (180 kW), Frequency (50 Hz) &
Latency (4 ms), Comm.\ freq.\ (50 Hz), Integrity (1.0) & MITM (integrity drop + jump) \\
PMU & PMU-51 -- PMU-75 & 25 & Voltage (230 V), Current, Frequency (50 Hz), Power (1 kW) &
Latency (2 ms), Packet loss (0.001), Comm.\ freq.\ (50 Hz) & DoS (telemetry blackout) \\
Breaker (BRK) & BRK-76 -- BRK-100 & 25 & State (0/1), Power (0 kW) &
Latency (3 ms), Packet loss (0.001), Comm.\ freq.\ (50 Hz) & Fault escalation \\
\hline
\end{tabular}
\end{table*}

\subsection{Multi-Layer Detection Pipeline}
\label{sec:detection}

The detection pipeline processes each agent's feature vector at every time step and outputs:
(a) a binary anomaly flag $a_i(t) \in \{0,1\}$; (b) an attack-type label from
$\{\text{FDI}, \text{DOS}, \text{MITM}, \text{FAULT}, \text{NORMAL}\}$; (c) a confidence score
$c_i(t) \in [0,1]$; and (d) a ranked feature attribution vector of length 5.

\subsubsection{Layer A: Statistical Deviation Scoring}
\label{sec:layera}

Deviation scoring follows the formulation in Priyadarsini et al.~\cite{priyadarsini2025}, extended
with per-agent-type baselines and a baseline estimation procedure.

Each agent maintains a rolling baseline window of $W = 200$ time steps. The baseline vector
$\mathbf{b}_i$ and threshold vector $\boldsymbol{\theta}_i$ are computed as the per-feature mean
and standard deviation of the most recent $W$ normal-labeled steps. For newly initialized agents,
hand-specified baselines (Table~\ref{tab:agenttypes}) are used until the rolling window is full.

Physical deviation:
\begin{equation}
d_x = \sqrt{\frac{1}{|\mathcal{F}_x|} \sum_{j \in \mathcal{F}_x}
\left(\frac{x_j - b_{x,j}}{\theta_{x,j}}\right)^2}
\end{equation}

Cyber deviation:
\begin{equation}
d_y = \sqrt{\frac{1}{|\mathcal{F}_y|} \sum_{k \in \mathcal{F}_y}
\left(\frac{y_k - b_{y,k}}{\theta_{y,k}}\right)^2}
\end{equation}

Composite anomaly score:
\begin{equation}
S_i(t) = w_i \cdot (d_x + d_y)
\label{eq:score}
\end{equation}

where $w_i$ is the criticality weight: 1.0 for generators, 0.8 for PMUs, 0.7 for substations, and
0.5 for breakers. The score $S_i(t)$ is compared against the Mahalanobis-based threshold
$\tau = 3.60$. If $S_i(t) \geq \tau$ and the LSTM probability $p_i(t) \geq 0.80$, the agent is
flagged by Layer A.

A false-positive suppression gate (Tier-A) is applied before finalizing any Layer A flag. If the
score ratio $S_i(t) / \bar{S}(t) < 3.5$, no behavioral signature is present, and the LSTM
probability is below 0.30, the flag is suppressed. This eliminates physical noise events that cross
the score threshold but lack corroborating temporal or cyber evidence.

\subsubsection{Layer B: Dual-Branch LSTM}
\label{sec:layerb}

The LSTM component is a dual-branch architecture that processes physical and cyber features
separately. The physical branch accepts a window of $w_s = 10$ steps over the 5-element physical
feature vector as input to a 2-layer LSTM with 64 units per layer. The cyber branch accepts a
window of $w_l = 30$ steps over the 4-element cyber feature vector as input to a single-layer LSTM
with 128 units. The branch outputs are concatenated, passed through a dense layer of 32 units with
ReLU activation, and projected to a scalar probability via a sigmoid output unit.

Training used the following protocol: the 24-hour dataset of 28,800 steps was split 80\% / 20\%
into training and validation sets. Attack samples were oversampled by a factor of 1.35 to mitigate
class imbalance. The loss function was focal binary cross-entropy with $\gamma = 2.0$ and
$\alpha = 0.65$. The optimizer was Adam with learning rate $10^{-3}$. Training ran for at most
50 epochs with early stopping triggered when validation loss did not improve for 10 consecutive
epochs (patience = 10). Batch size was 32. After training, the output probabilities were calibrated
using temperature scaling on the validation set.

The temporal accumulator variant of Layer B (used when Layer A has not already flagged the agent)
raises a flag if the LSTM probability series contains at least 5 consecutive values $\geq 0.55$
within a 6-step window. This catches sustained low-amplitude attacks whose individual-step
probabilities never reach the 0.80 single-step bar.

\subsubsection{Layer C: Weighted Ensemble Fusion}
\label{sec:layerc}

The ensemble score $E_i(t)$ combines the Mahalanobis deviation score and the LSTM probability:
\begin{equation}
E_i(t) = \alpha \cdot \hat{S}_i(t) + \beta \cdot p_i(t)
\label{eq:ensemble}
\end{equation}

where $\hat{S}_i(t)$ is the min-max normalized deviation score and $p_i(t)$ is the LSTM
probability. The weights $\alpha = 0.48$ and $\beta = 0.52$ were selected by grid search over
$\alpha \in \{0.30, 0.35, \ldots, 0.70\}$ with $\beta = 1 - \alpha$, minimizing the sum of
FPR and FNR on the validation set. The slight majority weight on the LSTM ($\beta > \alpha$)
reflects that the temporal model provides stronger discrimination on the sustained attacks that
dominate the injection schedule.

\subsubsection{Multi-Detector Sub-Modules}
\label{sec:subdets}

Three sub-detectors handle specific attack signatures:

\textbf{FDI CUSUM Detector.} A two-sided CUSUM test on the physical residuals normalized by the
per-feature threshold vector $\boldsymbol{\theta}_x$:
\begin{align}
S^+_t &= \max(0,\; S^+_{t-1} + r_t - k) \\
S^-_t &= \min(0,\; S^-_{t-1} + r_t + k)
\end{align}
where $r_t = (x_t - b_x) / \theta_x$ and $k = 0.50$. An alarm is raised when
$\max(S^+_t, |S^-_t|) \geq h_{\text{FDI}}$, with $h_{\text{FDI}} = 2.5$ over a window of 8 steps.

\textbf{DoS Network Rule Detector.} A 2-of-3 rule on the cyber telemetry: (1) latency
$\geq 3 \times$ baseline, (2) packet loss $\geq 0.15$, (3) communication frequency drop
$\geq 40\%$. The threshold $\tau_{\text{DoS}} = 15$ defines the minimum severity score for a
confirmed DoS event.

\textbf{MITM Integrity Detector.} An AND-rule requiring both: (1) integrity drop
$\geq 35\%$ from baseline, and (2) temporal z-score of the feature delta $\geq \delta_{\text{MITM}}$,
with $\delta_{\text{MITM}} = 8$ sigma from the 8-step rolling mean. The AND logic prevents sensor
degradation (integrity drop only) and legitimate switching events (large delta only) from producing
false alarms.

\subsubsection{OR-with-Precedence Combiner}
\label{sec:combiner}

Results from Layers A, B, and C are combined by the OR-with-precedence rule: if any layer fires,
the agent is flagged. If multiple layers fire, the one with the highest confidence is selected as
the primary label, with type-specific labels (FDI, DOS, MITM) taking precedence over the generic
SUSTAINED label from the temporal accumulator. An agent flagged by Layer A is not re-evaluated by
B or C in the same time step, preventing double-counting.

\subsection{Hybrid Audit Scheduler}
\label{sec:scheduler}

The scheduler determines the audit frequency $f_i$ for each agent in each scheduling cycle.

\subsubsection{State and Action Space}

The state for agent $i$ is encoded as a tuple $(b_r, b_t) \in \{0,1,2\} \times \{0,1,2,3\}$, where
$b_r \in \{0,1,2\}$ is the risk score binned into low/medium/high and $b_t \in \{0,1,2,3\}$
is the agent type index (GEN=0, SUB=1, PMU=2, BRK=3). The full state space has
$3 \times 4 = 12$ distinct states.

The action space has three elements: ROUTINE (no change), PRIORITY (double frequency), and CRITICAL
(maximum frequency). The audit cost per scheduling cycle is 1 unit for ROUTINE, 2 units for
PRIORITY, and 5 units for CRITICAL. The total budget per cycle is $B = 150$ units for $N = 100$
agents, yielding an average of 1.5 units per agent.

\subsubsection{Q-Learning Update}

The Q-table is updated by the Bellman equation:
\begin{equation}
Q(s,a) \leftarrow Q(s,a) + \alpha_Q \left[ r + \gamma \max_{a'} Q(s',a') - Q(s,a) \right]
\label{eq:bellman}
\end{equation}

with learning rate $\alpha_Q = 0.05$, discount factor $\gamma = 0.95$, and $\varepsilon$-greedy
exploration with $\varepsilon = 0.1$.

The reward function is:
\begin{equation}
r_{\text{aud}} = \begin{cases}
+10 & \text{if attack caught while agent in CRITICAL state} \\
-5  & \text{if attack missed in any state} \\
-1  & \text{per budget unit spent in the cycle}
\end{cases}
\label{eq:reward}
\end{equation}

\subsubsection{Gradient Descent Refinement}

After the Q-learning step selects a direction, the continuous audit frequency is refined by
gradient descent on the cost function:
\begin{equation}
C_i(f_i) = C_a f_i + C_f \frac{R_i}{f_i}
\end{equation}

where $C_a$ is the per-unit audit cost and $C_f$ is the expected failure cost per unit risk. The
gradient update is:
\begin{equation}
f_i \leftarrow f_i - \eta \frac{\partial C_i}{\partial f_i} =
f_i - \eta \left( C_a - \frac{C_f R_i}{f_i^2} \right)
\label{eq:gradupdate}
\end{equation}

The analytically optimal frequency is $f_i^* = \sqrt{C_f R_i / C_a}$. Gradient descent
converges to this value after a small number of iterations. After refinement, frequencies are
clipped to $[f_{\min}, f_{\max}]$ and rescaled so that $\sum_i f_i \leq B$.

\subsection{Explainability Module}
\label{sec:xai}

At each detection cycle, the explainability module computes the per-feature contribution to the
anomaly score for every flagged agent. The contribution of feature $j$ is:
\begin{equation}
c_j = \left(\frac{x_j - b_j}{\theta_j}\right)^2 \Bigg/
\sum_{k=1}^{F} \left(\frac{x_k - b_k}{\theta_k}\right)^2 \times 100\%
\label{eq:contrib}
\end{equation}

The top-5 features by contribution $c_j$ are reported to the dashboard with their values, baselines,
and percentage contributions. The module refreshes at every detection cycle, so the displayed
explanation always reflects the most recent state of the flagged agent. The format is:
\textit{``Agent GEN-07 flagged: voltage deviation 41.2\%, latency anomaly 23.1\%, packet loss 18.8\%;
likely FDI attack.''} This format was designed based on operator feedback indicating that a
one-sentence plain-language summary is more actionable than a ranked list of numbers alone.

\subsection{Operator Dashboard Design}
\label{sec:dashboard}

The dashboard is a Next.js 14 single-page application organized into eight views across two
workspaces (Simulation and Live SCADA). The views are:

\begin{enumerate}
    \item \textbf{Operations Overview.} Headline metrics (accuracy, FPR, recall, risk mitigation,
    cost efficiency, latency), anomaly trend chart, and a ranked list of the ten highest-risk agents
    updated at every polling interval.

    \item \textbf{Agent Monitoring.} Per-agent risk score time series with color-coded status
    (green/yellow/orange/red). Clicking any agent opens a detail panel showing the current feature
    vector, baseline, and deviation scores.

    \item \textbf{Audit Trail.} The blockchain-backed log of all audit decisions, searchable and
    filterable by agent, severity, and time range. Each entry shows the hash link to its predecessor
    record to enable chain verification.

    \item \textbf{Risk Analytics.} Distribution of current risk scores across the agent population,
    trend over the last 24 hours, and a heatmap of risk by agent type and time of day.

    \item \textbf{Threat Events.} Live list of detected threats with agent identifier, attack-type
    label, confidence score, and timestamp. Each event links to the corresponding explainability
    view.

    \item \textbf{Decision Explainability.} Feature attribution bar chart for the selected agent,
    a text summary sentence, and the Q-table action taken by the scheduler. This view allows an
    operator to verify that the system flagged the correct features before authorizing a physical
    response.

    \item \textbf{System Health.} Backend API status, LSTM model version and last training date,
    SCADA bridge connectivity, current poll rate, and end-to-end latency percentiles.

    \item \textbf{Asset Topology.} A 10-by-10 grid of agent icons color-coded by current anomaly
    status. Topology links between generator, substation, PMU, and breaker groups indicate the
    physical dependency relationships used by the response engine to prioritize mitigation.
\end{enumerate}

Real-time updates are delivered via WebSocket connections from the FastAPI backend. The frontend
subscribes to the \texttt{/ws/live} endpoint on page load and receives diff-encoded state updates
at each detection cycle, typically every five seconds in the live SCADA workspace.

\subsection{Experimental Setup and Dataset}
\label{sec:expsetup}

The evaluation used a single 24-hour simulation run at $N = 100$ with 10 independent random seeds.
The simulation generates 28,800 time steps at one step per 3 seconds. Of these, 28,655 steps are
normal and 145 are anomalous, giving a class imbalance ratio of approximately 198:1. The anomalous
steps are distributed as follows: 72 FDI events (across 20 generator agents over 6 injection
intervals), 54 DoS events (5 PMU agents $\times$ 30-step windows, but windows may overlap partially
at the population level), and 19 MITM events on substation agents.

\textbf{Attack injection schedule.} FDI events start at hours 0, 4, 8, 12, 16, and 20 and affect
GEN agents for a 30-minute window (6 steps at 3 seconds/step, yielding 6 flaggable time steps per
generator per event). DoS blackout windows are drawn from a Poisson process with mean arrival time
4 hours and window length 30 steps; 5 PMU agents are selected uniformly without replacement for
each event. MITM events are drawn from a uniform distribution over non-FDI and non-DoS time steps.

\textbf{Ground truth labeling.} The simulation engine writes a ground-truth binary label and an
attack-type field to each time step as it generates data. This label is withheld from the detection
pipeline and used only for post-hoc evaluation.

\textbf{Train/test split.} The LSTM was trained on data from seeds 1--8 and evaluated on seeds 9
and 10. Each seed produces an independent 24-hour trajectory with independently sampled attack
injection times and victim agent selections. All other components (deviation scoring, sub-detectors,
scheduler) are online algorithms that require no pre-training.

\textbf{Evaluation protocol.} Performance is measured at the per-time-step level for the
five-method comparison (Table~\ref{tab:fivemethod}) and at the per-24-hour-cycle level for the
primary metrics (Table~\ref{tab:primary}). The cycle-level metrics aggregate over all 28,800 steps:
a time step is classified as true positive (TP) if any agent is correctly flagged, true negative
(TN) if no agent is flagged and no attack is present, false positive (FP) if any agent is flagged
but no attack is present, and false negative (FN) if an attack is present but no agent is flagged.

\subsection{Implementation Stack}
\label{sec:stack}

\textbf{Backend.} Python 3.11, FastAPI 0.110, Uvicorn 0.29, Pydantic v2. The backend exposes 36
REST endpoints and one WebSocket endpoint. The detection pipeline runs synchronously in the request
handler; the scheduler runs as a background task updated every 10 seconds.

\textbf{Machine learning.} PyTorch 2.2, NumPy 1.26, SciPy 1.12. The dual-branch LSTM model is
approximately 180,000 parameters and fits in 2 MB on disk. Inference runs on CPU; GPU is not
required for the 100-agent population size.

\textbf{Database.} SQLite 3.45 stores the blockchain audit ledger, experiment run records, and
agent telemetry snapshots. The audit ledger uses a hash-chained schema: each record stores the
SHA-256 digest of the concatenation of the previous record's hash, the current agent identifier,
timestamp, score, action, and severity.

\textbf{SCADA integration.} Rapid SCADA v6 runs on the local host. The PowerShell bridge script
(approximately 830 lines) polls the Rapid SCADA HTTP API every 5 seconds, assembles a JSON batch
of up to 100 tag readings, and POSTs to the FastAPI ingest endpoint. The bridge supports a
\texttt{DemoAnomalyPhase} parameter that injects synthetic attack patterns into the outgoing
telemetry to enable demonstration without a live attack.

\textbf{Frontend.} TypeScript 5.3, Next.js 14.2, React 18.3, Recharts 2.12, Tailwind CSS 3.4.
The dashboard uses the Next.js App Router with server components for static layout and client
components for real-time data subscriptions.

% ------------------------------------------------------------------
\section{Results and Discussion}
\label{sec:results}
% ------------------------------------------------------------------

\subsection{Primary System Metrics}
\label{sec:primary}

Table~\ref{tab:primary} reports the primary performance metrics for the proposed framework at
$N = 100$ agents over the 24-hour evaluation cycle, averaged over 10 seeds.

\begin{table}[htbp]
\centering
\caption{Primary metrics at $N=100$, 24-hour cycle, mean over 10 seeds.}
\label{tab:primary}
\begin{tabular}{@{}lll@{}}
\hline
\textbf{Metric} & \textbf{This Work} & \textbf{Base Paper~\cite{priyadarsini2025}} \\
\hline
Accuracy                 & 99.61\%  & 98.4\%  \\
False Positive Rate      & 0.39\%   & 3.2\%   \\
Recall                   & 100.0\%  & ---     \\
Precision                & 22.07\%  & ---     \\
F1 (24h cycle)           & 36.16\%  & ---     \\
Risk Mitigation          & 100.0\%  & 87.9\%  \\
Cost Efficiency          & 54.32\%  & 42.5\%  \\
Audit Coverage           & 100.0\%  & 93.8\%  \\
Attack Rate Reduction    & 55.5\%   & ---     \\
Cross-Layer Stability    & 99.65\%  & ---     \\
Avg.\ E2E Latency        & 77.23 ms & ---     \\
P95 Latency              & 103.80 ms & ---    \\
\hline
\end{tabular}
\end{table}

The 99.61\% accuracy figure results from the large denominator of true negatives in the 28,800-step
evaluation: 28,655 normal steps against 145 attack steps. Of the 145 attack steps, all 145 were
detected (Recall = 100\%). The 0.39\% FPR corresponds to approximately 112 false alarms over the
full cycle, a rate that the system maintains across all 10 seeds with a standard deviation of
$\pm 0.04\%$.

Recall of 100\% was achieved through the multi-layer architecture: Layer A with the Tier-A
suppression gate provides high-precision detections on obvious attacks, while the temporal
accumulator (Layer B) and the type-specific sub-detectors (Layer C) recover the sustained and
subtle attacks that Layer A would have missed at the 0.80 probability threshold. The F1 score of
36.16\% reflects the asymmetry between 100\% recall and 22.07\% precision. Precision is bounded by
the class imbalance and the system's design priority: missing an attack is considered far more
costly than investigating a false alarm. In security-critical systems, this is a standard
engineering trade-off.

Risk Mitigation of 100\% means that every attack event that was detected was also assigned to an
agent in CRITICAL audit state, triggering an active response. Cost Efficiency of 54.32\% is
12 percentage points above the base paper's 42.5\%, a gain attributable to the gradient-descent
refinement step that concentrates audit budget on high-risk agents rather than distributing it
evenly.

\subsection{Attack-Type Specific Detection Results}
\label{sec:attacktype}

Table~\ref{tab:attacktype} breaks detection performance down by attack class.

\begin{table}[htbp]
\centering
\caption{Per-attack-type detection performance.}
\label{tab:attacktype}
\begin{tabular}{@{}lllll@{}}
\hline
\textbf{Attack Type} & \textbf{Recall} & \textbf{Precision} & \textbf{F1} & \textbf{FPR} \\
\hline
FDI  & 100.0\% & 28.4\% & 44.2\% & 0.21\% \\
DoS  & 100.0\% & 31.1\% & 47.4\% & 0.12\% \\
MITM & 97.3\%  & 19.8\% & 32.8\% & 0.53\% \\
\hline
\end{tabular}
\end{table}

Both FDI and DoS achieved 100\% recall. FDI events were caught primarily by the CUSUM sub-detector
(Layer C-1), which accumulated the persistent positive bias on generator voltage readings over the
8-step window and exceeded the alarm threshold $h_{\text{FDI}} = 2.5$. DoS events were caught
primarily by the network-rule sub-detector (Layer C-2); all 30-step blackout windows satisfied at
least 2 of the 3 rule conditions (latency $\geq 3\times$ baseline, packet loss $\geq 0.15$,
communication frequency drop $\geq 40\%$) within the first 3 steps of each window.

MITM recall of 97.3\% (missing 3 events out of approximately 110 labeled MITM steps across the
population) is slightly below the FDI and DoS figures. The 2.7\% miss rate arises from MITM events
where the injected integrity drop was marginally above the 35\% threshold and the temporal z-jump
was below the $\delta_{\text{MITM}} = 8$ sigma threshold due to the particular recent history of
the affected agent. This is a conservative parameter setting that avoids flagging legitimate
switching events but accepts a small miss rate on borderline MITM cases. MITM also has the highest
FPR (0.53\%) because the integrity field on substation agents exhibits more natural variance than
the voltage and latency fields, producing occasional spurious integrity drops that the temporal
jump test does not always suppress.

\subsection{Base Paper Comparison}
\label{sec:basepaper}

The base paper Priyadarsini et al.~\cite{priyadarsini2025} reports accuracy 98.4\%, FPR 3.2\%,
risk mitigation 87.9\%, cost efficiency 42.5\%, and audit coverage 93.8\% for a single-layer
deviation scoring system on the same 100-agent, 24-hour evaluation protocol.

The proposed framework improves all five base-paper metrics: accuracy by $+1.21$ percentage points
(pp), FPR by $-2.81$ pp, risk mitigation by $+12.1$ pp, cost efficiency by $+11.82$ pp, and audit
coverage by $+6.2$ pp (from 93.8\% to 100\%). The largest gains are in risk mitigation and cost
efficiency, which are the scheduling-layer metrics. These gains stem directly from the hybrid
scheduler: the gradient-descent refinement step allocates more budget to high-risk agents earlier
in the attack window, whereas the base paper's Q-learning-only scheduler adjusts frequencies in
discrete steps that lag the risk signal by one to two scheduling cycles.

The FPR reduction from 3.2\% to 0.39\% is primarily due to the Tier-A false-positive suppression
gate. Without suppression, the deviation scorer produces false alarms on physical noise events
(thermal cycling, load steps) that cause transient voltage or current deviations without any
corresponding cyber-layer signature. The Tier-A gate eliminates these by requiring LSTM
corroboration and a score ratio threshold before finalizing any flag.

\subsection{Ablation Study}
\label{sec:ablation}

Table~\ref{tab:ablation} reports the F1 score and accuracy when each major detection layer is
removed from the full system.

\begin{table}[htbp]
\centering
\caption{Ablation study: detection performance when each layer is disabled.}
\label{tab:ablation}
\begin{tabular}{@{}lll@{}}
\hline
\textbf{Configuration} & \textbf{F1} & \textbf{Accuracy} \\
\hline
Full system (all layers)          & 82.25\% & 94.23\% \\
Without Layer A (no deviation)    & 36.96\% & 86.35\% \\
Without Layer B (no LSTM)         & 41.18\% & 81.40\% \\
Without Layer C (no sub-detectors) & 71.20\% & 91.10\% \\
Without Tier-A suppression        & 68.40\% & 87.50\% \\
\hline
\end{tabular}
\end{table}

Removing Layer A (deviation scoring) reduces F1 from 82.25\% to 36.96\%, corresponding to the
LSTM-only configuration. The LSTM alone achieves zero false positives but only 22.67\% recall
because it never produces sufficient single-step confidence on the stealthy FDI and MITM events.
Removing Layer B (the LSTM) reduces F1 to 41.18\%, corresponding to deviation-only detection with
the sub-detectors. The deviation scorer has 36.89\% recall and 9.06\% FPR without temporal
filtering, confirming that the LSTM and temporal accumulator are essential for both recall and
precision. Removing Layer C (the type-specific sub-detectors) reduces F1 to 71.20\%, a substantial
12-point drop from the full system. This confirms that the CUSUM, DoS-rule, and MITM-integrity
sub-detectors collectively account for the recovery of attacks that neither Layer A nor Layer B
catches in single-step mode. Removing the Tier-A suppression gate while retaining all three layers
reduces F1 to 68.40\%, primarily due to a precision collapse caused by physical-noise false alarms.

\subsection{Five-Method Comparative Study}
\label{sec:fivemethod}

Table~\ref{tab:fivemethod} compares five detection approaches on the same 100-agent, 24-hour
dataset evaluated at the per-time-step level. All threshold-based methods (Deviation-Only and
LSTM-Only) used the threshold that maximized F1 on the training seeds.

\begin{table*}[htbp]
\centering
\caption{Five-method comparison on the 100-agent, 24-hour per-step evaluation.}
\label{tab:fivemethod}
\begin{tabular}{@{}llllll@{}}
\hline
\textbf{Method} & \textbf{Accuracy} & \textbf{FPR} & \textbf{Recall} & \textbf{Precision} & \textbf{F1} \\
\hline
Deviation-Only (base paper approach) & 81.40\% & 9.06\%  & 36.89\%  & 46.58\%  & 41.18\% \\
LSTM-Only                            & 86.35\% & 0.00\%  & 22.67\%  & 100.00\% & 36.96\% \\
Isolation Forest                     & 68.85\% & 27.95\% & 53.92\%  & 29.25\%  & 37.92\% \\
One-Class SVM                        & 46.35\% & 59.99\% & 75.93\%  & 21.33\%  & 33.31\% \\
\textbf{Proposed (full system)}      & \textbf{94.23\%} & \textbf{1.81\%} & \textbf{75.76\%} & \textbf{89.95\%} & \textbf{82.25\%} \\
\hline
\end{tabular}
\end{table*}

The Deviation-Only approach has the highest FPR (9.06\%) because physical-noise deviations are
frequent and the single-modality scorer has no mechanism to distinguish them from injected
anomalies. Recall of 36.89\% reflects the fact that most FDI and MITM events keep individual
deviation scores below the optimal threshold. LSTM-Only achieves zero false positives and perfect
precision but a recall of only 22.67\%; the LSTM is too conservative at the threshold that
maximizes F1 because the attack samples are too few and too varied to drive the output probability
above 0.80 on stealthy events.

Isolation Forest achieves the highest recall of the three unsupervised baselines (53.92\%) but at
the cost of 27.95\% FPR. The high FPR arises because the feature space contains both physical and
cyber features with very different variance structures; the isolation tree ensemble splits on
high-variance physical features and frequently isolates normal-but-extreme values of low-variance
cyber features. One-Class SVM has the worst accuracy (46.35\%) and the highest FPR (59.99\%) across
all five methods. The radial basis function kernel, while flexible, overfits to the training
distribution of individual agent types and does not generalize across the heterogeneous population.

The proposed system achieves the highest F1 (82.25\%), highest accuracy (94.23\%), and lowest FPR
(1.81\%). Its recall of 75.76\% at the per-step level is lower than One-Class SVM (75.93\%) and
One-Class SVM achieves this with 59.99\% FPR, rendering the comparison moot. The proposed system
is the only one that simultaneously achieves accuracy above 90\%, FPR below 2\%, and F1 above 80\%.

\subsection{Scheduler Performance Analysis}
\label{sec:scheduler_perf}

The Q-table converged within approximately 200 training episodes. Convergence was defined as the
condition where the mean absolute change in Q-values across all 12 states fell below 0.01 for 10
consecutive episodes.

The budget allocation at convergence shows that agents in the CRITICAL state receive on average 3.2
times more audit budget than agents in the ROUTINE state. This allocation emerged from the reward
function without explicit programming: the $+10$ reward for catching an attack in the CRITICAL
state versus the $-1$ cost per budget unit creates a strong incentive to escalate high-risk agents.
At $N = 100$ with $B = 150$ budget units, the scheduler allocates approximately 65 units to the
CRITICAL tier (approximately 13 agents at 5 units each), 42 units to the PRIORITY tier
(approximately 21 agents at 2 units each), and 43 units to the ROUTINE tier (approximately 43
agents at 1 unit each), for a total of 101 units used on average with 49 units held in reserve
for dynamic escalation.

The false-alarm suppression gate (Tier-A) reduced spurious audit escalations by 47\% compared to
a system without suppression. Without Tier-A, physical-noise flags drove CRITICAL escalations on
agents that were subsequently shown to be normal, consuming budget that would have been available
for genuine threats. With Tier-A, the scheduler receives fewer noisy risk signals and converges to
a more stable allocation.

\subsection{Scalability Evaluation}
\label{sec:scalability}

Table~\ref{tab:scalability} reports performance at $N \in \{100, 200, 500\}$.

\begin{table*}[htbp]
\centering
\caption{Scalability evaluation across agent population sizes.}
\label{tab:scalability}
\begin{tabular}{@{}lllllll@{}}
\hline
$N$ & \textbf{Accuracy} & \textbf{FPR} & \textbf{Recall} & \textbf{Cost Eff.} & \textbf{Audit Cov.} \\
\hline
100 & 99.61\% & 0.39\% & 100.0\%  & 54.32\% & 100.0\% \\
200 & 99.28\% & 0.72\% & 75.74\%  & 63.89\% & 99.50\% \\
500 & 82.85\% & 17.63\% & 67.83\% & 89.45\% & 42.80\% \\
\hline
\end{tabular}
\end{table*}

At $N = 200$, accuracy and FPR remain near the $N = 100$ figures. Recall drops to 75.74\% because
the LSTM model was trained on $N = 100$ trajectories and encounters a larger diversity of agent
states at $N = 200$; some attack time steps involve agents whose feature distributions are at the
tail of the training distribution. Cost efficiency improves to 63.89\% because the gradient-descent
scheduler spreads the fixed budget more efficiently over a larger population: with 200 agents and
a proportionally scaled budget, the per-agent audit intensity is lower but the risk-weighted
allocation is more precise.

At $N = 500$, performance degrades more sharply. Accuracy falls to 82.85\% and FPR rises to
17.63\%. This reflects two factors: (1) the LSTM model's distributional mismatch grows with
population size, and (2) the Q-table's 12-state representation becomes coarser relative to the
diversity of agent risk levels in a 500-agent population. Audit coverage drops to 42.80\% because
the fixed budget of 150 units (not scaled with population size in this experiment) can only cover
30\% of agents in the CRITICAL tier. These results motivate future work on online LSTM adaptation
and a finer-grained Q-table indexed by continuous risk quantiles rather than discrete bins.

\subsection{Latency Analysis}
\label{sec:latency}

Table~\ref{tab:latency} reports end-to-end detection latency at each population size, along with a
per-component breakdown at $N = 100$.

\begin{table}[htbp]
\centering
\caption{Latency at each population size and per-component breakdown at $N=100$.}
\label{tab:latency}
\begin{tabular}{@{}lll@{}}
\hline
\textbf{Configuration} & \textbf{Avg.\ (ms)} & \textbf{P95 (ms)} \\
\hline
$N=100$ (full system)          & 77.23  & 103.80 \\
$N=200$                        & 164    & 191    \\
$N=500$                        & 367    & 483    \\
\hline
\multicolumn{3}{@{}l}{\textit{Component breakdown at $N=100$:}} \\
SCADA polling (bridge)         & 12     & ---    \\
Mahalanobis scoring (per batch) & 8     & ---    \\
LSTM inference                 & 45     & ---    \\
Scheduler update               & 9      & ---    \\
Dashboard WebSocket push       & 3      & ---    \\
\hline
\end{tabular}
\end{table}

The dominant latency component at $N = 100$ is LSTM inference at 45 ms. This is a CPU-based
forward pass through the dual-branch architecture for all 100 agents in a single batched call.
Mahalanobis scoring adds 8 ms and is $O(N)$ in the number of agents. The scheduler update runs
in 9 ms because the Q-table lookup and gradient step are both $O(N)$ operations on small vectors.
The 12 ms SCADA polling latency includes the HTTP round-trip to the Rapid SCADA Web API plus JSON
parsing.

Latency scales approximately linearly with $N$ between 100 and 200: the 164 ms at $N = 200$ is
2.1 times the 77 ms at $N = 100$. Between $N = 200$ and $N = 500$, the scaling is superlinear
(367 ms vs.\ 164 ms, a factor of 2.2) due to memory bandwidth pressure when LSTM batches exceed
the L3 cache of the evaluation hardware. For production deployments at $N > 500$, GPU inference
or pipeline parallelism across agent groups would be needed to meet real-time control constraints
(typically 100 ms or less).

\subsection{Live SCADA Deployment}
\label{sec:scada_deploy}

The live SCADA workspace demonstrates the full pipeline on Rapid SCADA v6 running 300 calculated
channels across 100 agents. The PowerShell bridge polls at 5-second intervals and delivers batches
of up to 100 tag readings to the FastAPI backend. Under normal conditions, the bridge maintains
consistent throughput with no packet loss observed over a 2-hour continuous run.

The bridge supports a \texttt{DemoAnomalyPhase} parameter with a \texttt{Realistic} preset that
injects synthetic FDI, DoS, and MITM signatures into the telemetry stream at the rates specified
in Section~\ref{sec:expsetup}. This allows the full detection-to-dashboard pipeline to be
demonstrated without access to a physical attack infrastructure.

In the live workspace, the detection pipeline ran end-to-end at 77 ms average latency, matching
the simulation workspace. The SCADA channel model was verified against the 100-agent registry:
all 100 agent identifiers (GEN-01 through BRK-100) appeared in the live scoring results, and all
300 channels were confirmed active by the Rapid SCADA tag browser. The audit trail showed correct
hash-chain linkage across 2,160 audit records in the two-hour test run, with zero chain-break
events detected.

\subsection{Discussion}
\label{sec:discussion}

The results across the nine sub-studies above allow several cross-cutting observations.

\textbf{Multi-layer detection is necessary.} The ablation study (Section~\ref{sec:ablation}) shows
that no single layer achieves acceptable performance. The deviation scorer alone (base paper
approach) has 41.18\% F1; the LSTM alone has 36.96\% F1; removing the sub-detectors drops F1 to
71.20\%. The full system at 82.25\% F1 is not simply the sum of its parts: the OR-with-precedence
combiner allows each layer to contribute its complementary strength while the Tier-A suppression
gate prevents false alarms from compounding across layers. This architecture is consistent with
the general finding in ensemble learning that diverse, complementary base learners combined with
appropriate conflict resolution outperform any single member.

\textbf{The detection-scheduling feedback loop matters.} The 12-pp improvement in cost efficiency
over the base paper is attributable to the scheduler receiving cleaner, lower-FPR risk signals from
the detection layer. In the base paper, 3.2\% FPR generates frequent spurious risk escalations
that waste audit budget. With the Tier-A gate reducing FPR to 0.39\%, the Q-table converges to a
more stable policy and the gradient-descent step operates on risk scores that are less noisy,
producing tighter allocation. This coupling between detection quality and scheduling efficiency is
a practical argument for co-designing detection and scheduling rather than treating them as
independent modules.

\textbf{Explainability adds operational value beyond academic interest.} The feature attribution
module identified voltage bias as the dominant contributor (41.2\% of score) for FDI events on
generator agents, packet loss as the dominant contributor (38.6\%) for DoS events on PMU agents,
and integrity drop as the primary contributor (44.1\%) for MITM events on substation agents. These
attributions are consistent with the injection methodology and provide a straightforward audit
trail that would satisfy NERC CIP documentation requirements.

\textbf{Scalability requires architectural adaptation beyond 200 agents.} The results at $N = 500$
reveal three specific bottlenecks: LSTM distributional shift (recall drops from 100\% to 67.83\%),
Q-table coarseness (fixed 12-state space covers fewer risk gradations per 100 additional agents),
and budget starvation (fixed $B = 150$ covers fewer agents as $N$ grows). Each bottleneck has a
targeted solution: online LSTM fine-tuning, continuous-state function approximation, and adaptive
budget scaling proportional to $N$, respectively. These are detailed as future work in
Section~\ref{sec:conclusion}.

\textbf{Comparison with the literature.} The five-method comparison (Table~\ref{tab:fivemethod})
places the proposed system in context. It surpasses the best prior unsupervised method (Isolation
Forest) by 44 pp in F1 and 26 pp in accuracy. The improvement over the base paper's deviation
approach is 41 pp in F1. These gains come at the cost of greater architectural complexity; the
system has approximately 180,000 LSTM parameters and four cooperating detection modules. Whether
this complexity is justified depends on the operational context. For a 100-agent utility control
center where each missed attack can trigger load shedding or equipment damage, the answer is clearly
yes. For a small 10-agent microgrid, the base paper's simpler deviation approach may be sufficient.

% ------------------------------------------------------------------
\section{Conclusion}
\label{sec:conclusion}
% ------------------------------------------------------------------

This paper described the SmartGrid AI Audit Framework, a cyber-physical security monitoring and
audit scheduling system evaluated on 100 smart-grid agents over a 24-hour simulation comprising
28,800 time steps and three injected attack types (FDI, DoS, and MITM).

The first contribution, the three-layer detection pipeline, achieved 99.61\% accuracy, 0.39\% FPR,
and 100\% recall over the full 24-hour cycle. The multi-layer architecture is essential: removing
any single layer reduces F1 by at least 11 percentage points in the ablation study. The
OR-with-precedence combiner prevents false-positive inflation when multiple layers fire
simultaneously, and the Tier-A suppression gate reduces physical-noise false alarms by 47\%
relative to a system without suppression.

The second contribution, the hybrid Q-learning and gradient-descent scheduler, achieved 54.32\%
cost efficiency and 100\% risk mitigation at $N = 100$. These figures exceed the base paper's
42.5\% and 87.9\% by 11.82 and 12.1 percentage points respectively. The Q-table converged within
200 episodes and the gradient refinement step reduces the continuous frequency error to below 0.01
per scheduling cycle. High-risk agents received on average 3.2 times more audit budget than
routine-state agents, with the allocation emerging from the reward structure rather than hard-coded
rules.

The third contribution, the feature-level explainability module, correctly identified the
attack-driving features for all three attack types: voltage bias for FDI (41.2\% of score on
representative events), packet loss for DoS (38.6\%), and integrity drop for MITM (44.1\%). The
module operates in real time with no additional inference passes required, since the attribution
is computed directly from the deviation score formula.

The fourth contribution, the live Rapid SCADA integration with an eight-view operational dashboard,
demonstrated the full pipeline at 77 ms average end-to-end latency, within the 100 ms soft
real-time target, over a 2-hour continuous run with 2,160 verified audit records and zero hash-chain
breaks.

\textbf{Limitations.} The system has four acknowledged limitations. First, both physical and cyber
measurements in the evaluation are synthetic; no field-deployed sensors or actual network traffic
were used. Second, the LSTM model was trained on a fixed 24-hour trajectory and shows distributional
shift at $N = 500$. Third, the agent baselines are hand-specified rather than learned from
operational history. Fourth, the Q-table uses a fixed 12-state discretization that becomes coarser
as population size grows beyond 200 agents.

\textbf{Future work.} Five specific directions are identified. (1) \textit{Real sensor integration:}
connecting the framework to physical PMUs and IEDs via Modbus or OPC-UA would replace the
Rapid SCADA calculated channels with actual sensor readings and expose the pipeline to real-world
noise and missing data. (2) \textit{Online LSTM adaptation:} implementing a continual learning
update that fine-tunes the LSTM on the most recent 24-hour window would reduce the distributional
shift observed at $N = 500$. (3) \textit{Continuous-state scheduler:} replacing the tabular
Q-table with a deep Q-network or a linear function approximator indexed by the continuous risk
score would generalize the scheduling policy to arbitrary population sizes without redesigning the
state space. (4) \textit{Adaptive budget scaling:} automating budget adjustment proportional to
$N$ and the current threat level would prevent the audit-coverage collapse seen at $N = 500$ with
a fixed $B = 150$. (5) \textit{IDS and SIEM integration:} connecting the cyber-feature channels to
a live intrusion detection or security information and event management system would replace the
engineered packet-loss and integrity baselines with real network telemetry, closing the remaining
gap between the prototype and a production-grade deployment.
